%% kapa: exam=1 round=34 year=2023 period=1 form=A subject=economics subject_ko=경제학원론 seq=22 number=62
%% source: Q-Net 감정평가사 기출문제 (cst00309) artlSeq=5211789 file=raw/1008/5211789/2023년도 제34회 감정평가사 1차시험 1교시.pdf
%% conversion: kordoc 4.24.51 -> tools/gen1.py (auto); review=proofread-sonnet; flags=none
\begin{question}{62}
국내총생산에 관한 설명으로 옳지 않은 것은?
\begin{choices}
  \item 국내총생산은 시장에서 거래되는 최종생산물만 포함한다.
  \item 국내순생산은 국내총생산에서 고정자본소모를 제외한 부분이다.
  \item 명목국내총생산은 재화와 서비스의 생산의 가치를 경상가격으로 계산한 것이다.
  \item 3면 등가의 원칙으로 국내총생산은 국내총소득과 일치한다.
  \item 국내총생산은 요소비용국내소득에 순간접세와 고정자본소모를 더한 것이다.
\end{choices}
\end{question}
\answer{1}
